\documentclass{article} % For LaTeX2e
\usepackage{iclr2025_conference,times}

% Optional math commands from https://github.com/goodfeli/dlbook_notation.
%%%%% NEW MATH DEFINITIONS %%%%%

\usepackage{amsmath,amsfonts,bm}

% Mark sections of captions for referring to divisions of figures
\newcommand{\figleft}{{\em (Left)}}
\newcommand{\figcenter}{{\em (Center)}}
\newcommand{\figright}{{\em (Right)}}
\newcommand{\figtop}{{\em (Top)}}
\newcommand{\figbottom}{{\em (Bottom)}}
\newcommand{\captiona}{{\em (a)}}
\newcommand{\captionb}{{\em (b)}}
\newcommand{\captionc}{{\em (c)}}
\newcommand{\captiond}{{\em (d)}}

% Highlight a newly defined term
\newcommand{\newterm}[1]{{\bf #1}}


% Figure reference, lower-case.
\def\figref#1{figure~\ref{#1}}
% Figure reference, capital. For start of sentence
\def\Figref#1{Figure~\ref{#1}}
\def\twofigref#1#2{figures \ref{#1} and \ref{#2}}
\def\quadfigref#1#2#3#4{figures \ref{#1}, \ref{#2}, \ref{#3} and \ref{#4}}
% Section reference, lower-case.
\def\secref#1{section~\ref{#1}}
% Section reference, capital.
\def\Secref#1{Section~\ref{#1}}
% Reference to two sections.
\def\twosecrefs#1#2{sections \ref{#1} and \ref{#2}}
% Reference to three sections.
\def\secrefs#1#2#3{sections \ref{#1}, \ref{#2} and \ref{#3}}
% Reference to an equation, lower-case.
\def\eqref#1{equation~\ref{#1}}
% Reference to an equation, upper case
\def\Eqref#1{Equation~\ref{#1}}
% A raw reference to an equation---avoid using if possible
\def\plaineqref#1{\ref{#1}}
% Reference to a chapter, lower-case.
\def\chapref#1{chapter~\ref{#1}}
% Reference to an equation, upper case.
\def\Chapref#1{Chapter~\ref{#1}}
% Reference to a range of chapters
\def\rangechapref#1#2{chapters\ref{#1}--\ref{#2}}
% Reference to an algorithm, lower-case.
\def\algref#1{algorithm~\ref{#1}}
% Reference to an algorithm, upper case.
\def\Algref#1{Algorithm~\ref{#1}}
\def\twoalgref#1#2{algorithms \ref{#1} and \ref{#2}}
\def\Twoalgref#1#2{Algorithms \ref{#1} and \ref{#2}}
% Reference to a part, lower case
\def\partref#1{part~\ref{#1}}
% Reference to a part, upper case
\def\Partref#1{Part~\ref{#1}}
\def\twopartref#1#2{parts \ref{#1} and \ref{#2}}

\def\ceil#1{\lceil #1 \rceil}
\def\floor#1{\lfloor #1 \rfloor}
\def\1{\bm{1}}
\newcommand{\train}{\mathcal{D}}
\newcommand{\valid}{\mathcal{D_{\mathrm{valid}}}}
\newcommand{\test}{\mathcal{D_{\mathrm{test}}}}

\def\eps{{\epsilon}}


% Random variables
\def\reta{{\textnormal{$\eta$}}}
\def\ra{{\textnormal{a}}}
\def\rb{{\textnormal{b}}}
\def\rc{{\textnormal{c}}}
\def\rd{{\textnormal{d}}}
\def\re{{\textnormal{e}}}
\def\rf{{\textnormal{f}}}
\def\rg{{\textnormal{g}}}
\def\rh{{\textnormal{h}}}
\def\ri{{\textnormal{i}}}
\def\rj{{\textnormal{j}}}
\def\rk{{\textnormal{k}}}
\def\rl{{\textnormal{l}}}
% rm is already a command, just don't name any random variables m
\def\rn{{\textnormal{n}}}
\def\ro{{\textnormal{o}}}
\def\rp{{\textnormal{p}}}
\def\rq{{\textnormal{q}}}
\def\rr{{\textnormal{r}}}
\def\rs{{\textnormal{s}}}
\def\rt{{\textnormal{t}}}
\def\ru{{\textnormal{u}}}
\def\rv{{\textnormal{v}}}
\def\rw{{\textnormal{w}}}
\def\rx{{\textnormal{x}}}
\def\ry{{\textnormal{y}}}
\def\rz{{\textnormal{z}}}

% Random vectors
\def\rvepsilon{{\mathbf{\epsilon}}}
\def\rvtheta{{\mathbf{\theta}}}
\def\rva{{\mathbf{a}}}
\def\rvb{{\mathbf{b}}}
\def\rvc{{\mathbf{c}}}
\def\rvd{{\mathbf{d}}}
\def\rve{{\mathbf{e}}}
\def\rvf{{\mathbf{f}}}
\def\rvg{{\mathbf{g}}}
\def\rvh{{\mathbf{h}}}
\def\rvu{{\mathbf{i}}}
\def\rvj{{\mathbf{j}}}
\def\rvk{{\mathbf{k}}}
\def\rvl{{\mathbf{l}}}
\def\rvm{{\mathbf{m}}}
\def\rvn{{\mathbf{n}}}
\def\rvo{{\mathbf{o}}}
\def\rvp{{\mathbf{p}}}
\def\rvq{{\mathbf{q}}}
\def\rvr{{\mathbf{r}}}
\def\rvs{{\mathbf{s}}}
\def\rvt{{\mathbf{t}}}
\def\rvu{{\mathbf{u}}}
\def\rvv{{\mathbf{v}}}
\def\rvw{{\mathbf{w}}}
\def\rvx{{\mathbf{x}}}
\def\rvy{{\mathbf{y}}}
\def\rvz{{\mathbf{z}}}

% Elements of random vectors
\def\erva{{\textnormal{a}}}
\def\ervb{{\textnormal{b}}}
\def\ervc{{\textnormal{c}}}
\def\ervd{{\textnormal{d}}}
\def\erve{{\textnormal{e}}}
\def\ervf{{\textnormal{f}}}
\def\ervg{{\textnormal{g}}}
\def\ervh{{\textnormal{h}}}
\def\ervi{{\textnormal{i}}}
\def\ervj{{\textnormal{j}}}
\def\ervk{{\textnormal{k}}}
\def\ervl{{\textnormal{l}}}
\def\ervm{{\textnormal{m}}}
\def\ervn{{\textnormal{n}}}
\def\ervo{{\textnormal{o}}}
\def\ervp{{\textnormal{p}}}
\def\ervq{{\textnormal{q}}}
\def\ervr{{\textnormal{r}}}
\def\ervs{{\textnormal{s}}}
\def\ervt{{\textnormal{t}}}
\def\ervu{{\textnormal{u}}}
\def\ervv{{\textnormal{v}}}
\def\ervw{{\textnormal{w}}}
\def\ervx{{\textnormal{x}}}
\def\ervy{{\textnormal{y}}}
\def\ervz{{\textnormal{z}}}

% Random matrices
\def\rmA{{\mathbf{A}}}
\def\rmB{{\mathbf{B}}}
\def\rmC{{\mathbf{C}}}
\def\rmD{{\mathbf{D}}}
\def\rmE{{\mathbf{E}}}
\def\rmF{{\mathbf{F}}}
\def\rmG{{\mathbf{G}}}
\def\rmH{{\mathbf{H}}}
\def\rmI{{\mathbf{I}}}
\def\rmJ{{\mathbf{J}}}
\def\rmK{{\mathbf{K}}}
\def\rmL{{\mathbf{L}}}
\def\rmM{{\mathbf{M}}}
\def\rmN{{\mathbf{N}}}
\def\rmO{{\mathbf{O}}}
\def\rmP{{\mathbf{P}}}
\def\rmQ{{\mathbf{Q}}}
\def\rmR{{\mathbf{R}}}
\def\rmS{{\mathbf{S}}}
\def\rmT{{\mathbf{T}}}
\def\rmU{{\mathbf{U}}}
\def\rmV{{\mathbf{V}}}
\def\rmW{{\mathbf{W}}}
\def\rmX{{\mathbf{X}}}
\def\rmY{{\mathbf{Y}}}
\def\rmZ{{\mathbf{Z}}}

% Elements of random matrices
\def\ermA{{\textnormal{A}}}
\def\ermB{{\textnormal{B}}}
\def\ermC{{\textnormal{C}}}
\def\ermD{{\textnormal{D}}}
\def\ermE{{\textnormal{E}}}
\def\ermF{{\textnormal{F}}}
\def\ermG{{\textnormal{G}}}
\def\ermH{{\textnormal{H}}}
\def\ermI{{\textnormal{I}}}
\def\ermJ{{\textnormal{J}}}
\def\ermK{{\textnormal{K}}}
\def\ermL{{\textnormal{L}}}
\def\ermM{{\textnormal{M}}}
\def\ermN{{\textnormal{N}}}
\def\ermO{{\textnormal{O}}}
\def\ermP{{\textnormal{P}}}
\def\ermQ{{\textnormal{Q}}}
\def\ermR{{\textnormal{R}}}
\def\ermS{{\textnormal{S}}}
\def\ermT{{\textnormal{T}}}
\def\ermU{{\textnormal{U}}}
\def\ermV{{\textnormal{V}}}
\def\ermW{{\textnormal{W}}}
\def\ermX{{\textnormal{X}}}
\def\ermY{{\textnormal{Y}}}
\def\ermZ{{\textnormal{Z}}}

% Vectors
\def\vzero{{\bm{0}}}
\def\vone{{\bm{1}}}
\def\vmu{{\bm{\mu}}}
\def\vtheta{{\bm{\theta}}}
\def\va{{\bm{a}}}
\def\vb{{\bm{b}}}
\def\vc{{\bm{c}}}
\def\vd{{\bm{d}}}
\def\ve{{\bm{e}}}
\def\vf{{\bm{f}}}
\def\vg{{\bm{g}}}
\def\vh{{\bm{h}}}
\def\vi{{\bm{i}}}
\def\vj{{\bm{j}}}
\def\vk{{\bm{k}}}
\def\vl{{\bm{l}}}
\def\vm{{\bm{m}}}
\def\vn{{\bm{n}}}
\def\vo{{\bm{o}}}
\def\vp{{\bm{p}}}
\def\vq{{\bm{q}}}
\def\vr{{\bm{r}}}
\def\vs{{\bm{s}}}
\def\vt{{\bm{t}}}
\def\vu{{\bm{u}}}
\def\vv{{\bm{v}}}
\def\vw{{\bm{w}}}
\def\vx{{\bm{x}}}
\def\vy{{\bm{y}}}
\def\vz{{\bm{z}}}

% Elements of vectors
\def\evalpha{{\alpha}}
\def\evbeta{{\beta}}
\def\evepsilon{{\epsilon}}
\def\evlambda{{\lambda}}
\def\evomega{{\omega}}
\def\evmu{{\mu}}
\def\evpsi{{\psi}}
\def\evsigma{{\sigma}}
\def\evtheta{{\theta}}
\def\eva{{a}}
\def\evb{{b}}
\def\evc{{c}}
\def\evd{{d}}
\def\eve{{e}}
\def\evf{{f}}
\def\evg{{g}}
\def\evh{{h}}
\def\evi{{i}}
\def\evj{{j}}
\def\evk{{k}}
\def\evl{{l}}
\def\evm{{m}}
\def\evn{{n}}
\def\evo{{o}}
\def\evp{{p}}
\def\evq{{q}}
\def\evr{{r}}
\def\evs{{s}}
\def\evt{{t}}
\def\evu{{u}}
\def\evv{{v}}
\def\evw{{w}}
\def\evx{{x}}
\def\evy{{y}}
\def\evz{{z}}

% Matrix
\def\mA{{\bm{A}}}
\def\mB{{\bm{B}}}
\def\mC{{\bm{C}}}
\def\mD{{\bm{D}}}
\def\mE{{\bm{E}}}
\def\mF{{\bm{F}}}
\def\mG{{\bm{G}}}
\def\mH{{\bm{H}}}
\def\mI{{\bm{I}}}
\def\mJ{{\bm{J}}}
\def\mK{{\bm{K}}}
\def\mL{{\bm{L}}}
\def\mM{{\bm{M}}}
\def\mN{{\bm{N}}}
\def\mO{{\bm{O}}}
\def\mP{{\bm{P}}}
\def\mQ{{\bm{Q}}}
\def\mR{{\bm{R}}}
\def\mS{{\bm{S}}}
\def\mT{{\bm{T}}}
\def\mU{{\bm{U}}}
\def\mV{{\bm{V}}}
\def\mW{{\bm{W}}}
\def\mX{{\bm{X}}}
\def\mY{{\bm{Y}}}
\def\mZ{{\bm{Z}}}
\def\mBeta{{\bm{\beta}}}
\def\mPhi{{\bm{\Phi}}}
\def\mLambda{{\bm{\Lambda}}}
\def\mSigma{{\bm{\Sigma}}}

% Tensor
\DeclareMathAlphabet{\mathsfit}{\encodingdefault}{\sfdefault}{m}{sl}
\SetMathAlphabet{\mathsfit}{bold}{\encodingdefault}{\sfdefault}{bx}{n}
\newcommand{\tens}[1]{\bm{\mathsfit{#1}}}
\def\tA{{\tens{A}}}
\def\tB{{\tens{B}}}
\def\tC{{\tens{C}}}
\def\tD{{\tens{D}}}
\def\tE{{\tens{E}}}
\def\tF{{\tens{F}}}
\def\tG{{\tens{G}}}
\def\tH{{\tens{H}}}
\def\tI{{\tens{I}}}
\def\tJ{{\tens{J}}}
\def\tK{{\tens{K}}}
\def\tL{{\tens{L}}}
\def\tM{{\tens{M}}}
\def\tN{{\tens{N}}}
\def\tO{{\tens{O}}}
\def\tP{{\tens{P}}}
\def\tQ{{\tens{Q}}}
\def\tR{{\tens{R}}}
\def\tS{{\tens{S}}}
\def\tT{{\tens{T}}}
\def\tU{{\tens{U}}}
\def\tV{{\tens{V}}}
\def\tW{{\tens{W}}}
\def\tX{{\tens{X}}}
\def\tY{{\tens{Y}}}
\def\tZ{{\tens{Z}}}


% Graph
\def\gA{{\mathcal{A}}}
\def\gB{{\mathcal{B}}}
\def\gC{{\mathcal{C}}}
\def\gD{{\mathcal{D}}}
\def\gE{{\mathcal{E}}}
\def\gF{{\mathcal{F}}}
\def\gG{{\mathcal{G}}}
\def\gH{{\mathcal{H}}}
\def\gI{{\mathcal{I}}}
\def\gJ{{\mathcal{J}}}
\def\gK{{\mathcal{K}}}
\def\gL{{\mathcal{L}}}
\def\gM{{\mathcal{M}}}
\def\gN{{\mathcal{N}}}
\def\gO{{\mathcal{O}}}
\def\gP{{\mathcal{P}}}
\def\gQ{{\mathcal{Q}}}
\def\gR{{\mathcal{R}}}
\def\gS{{\mathcal{S}}}
\def\gT{{\mathcal{T}}}
\def\gU{{\mathcal{U}}}
\def\gV{{\mathcal{V}}}
\def\gW{{\mathcal{W}}}
\def\gX{{\mathcal{X}}}
\def\gY{{\mathcal{Y}}}
\def\gZ{{\mathcal{Z}}}

% Sets
\def\sA{{\mathbb{A}}}
\def\sB{{\mathbb{B}}}
\def\sC{{\mathbb{C}}}
\def\sD{{\mathbb{D}}}
% Don't use a set called E, because this would be the same as our symbol
% for expectation.
\def\sF{{\mathbb{F}}}
\def\sG{{\mathbb{G}}}
\def\sH{{\mathbb{H}}}
\def\sI{{\mathbb{I}}}
\def\sJ{{\mathbb{J}}}
\def\sK{{\mathbb{K}}}
\def\sL{{\mathbb{L}}}
\def\sM{{\mathbb{M}}}
\def\sN{{\mathbb{N}}}
\def\sO{{\mathbb{O}}}
\def\sP{{\mathbb{P}}}
\def\sQ{{\mathbb{Q}}}
\def\sR{{\mathbb{R}}}
\def\sS{{\mathbb{S}}}
\def\sT{{\mathbb{T}}}
\def\sU{{\mathbb{U}}}
\def\sV{{\mathbb{V}}}
\def\sW{{\mathbb{W}}}
\def\sX{{\mathbb{X}}}
\def\sY{{\mathbb{Y}}}
\def\sZ{{\mathbb{Z}}}

% Entries of a matrix
\def\emLambda{{\Lambda}}
\def\emA{{A}}
\def\emB{{B}}
\def\emC{{C}}
\def\emD{{D}}
\def\emE{{E}}
\def\emF{{F}}
\def\emG{{G}}
\def\emH{{H}}
\def\emI{{I}}
\def\emJ{{J}}
\def\emK{{K}}
\def\emL{{L}}
\def\emM{{M}}
\def\emN{{N}}
\def\emO{{O}}
\def\emP{{P}}
\def\emQ{{Q}}
\def\emR{{R}}
\def\emS{{S}}
\def\emT{{T}}
\def\emU{{U}}
\def\emV{{V}}
\def\emW{{W}}
\def\emX{{X}}
\def\emY{{Y}}
\def\emZ{{Z}}
\def\emSigma{{\Sigma}}

% entries of a tensor
% Same font as tensor, without \bm wrapper
\newcommand{\etens}[1]{\mathsfit{#1}}
\def\etLambda{{\etens{\Lambda}}}
\def\etA{{\etens{A}}}
\def\etB{{\etens{B}}}
\def\etC{{\etens{C}}}
\def\etD{{\etens{D}}}
\def\etE{{\etens{E}}}
\def\etF{{\etens{F}}}
\def\etG{{\etens{G}}}
\def\etH{{\etens{H}}}
\def\etI{{\etens{I}}}
\def\etJ{{\etens{J}}}
\def\etK{{\etens{K}}}
\def\etL{{\etens{L}}}
\def\etM{{\etens{M}}}
\def\etN{{\etens{N}}}
\def\etO{{\etens{O}}}
\def\etP{{\etens{P}}}
\def\etQ{{\etens{Q}}}
\def\etR{{\etens{R}}}
\def\etS{{\etens{S}}}
\def\etT{{\etens{T}}}
\def\etU{{\etens{U}}}
\def\etV{{\etens{V}}}
\def\etW{{\etens{W}}}
\def\etX{{\etens{X}}}
\def\etY{{\etens{Y}}}
\def\etZ{{\etens{Z}}}

% The true underlying data generating distribution
\newcommand{\pdata}{p_{\rm{data}}}
% The empirical distribution defined by the training set
\newcommand{\ptrain}{\hat{p}_{\rm{data}}}
\newcommand{\Ptrain}{\hat{P}_{\rm{data}}}
% The model distribution
\newcommand{\pmodel}{p_{\rm{model}}}
\newcommand{\Pmodel}{P_{\rm{model}}}
\newcommand{\ptildemodel}{\tilde{p}_{\rm{model}}}
% Stochastic autoencoder distributions
\newcommand{\pencode}{p_{\rm{encoder}}}
\newcommand{\pdecode}{p_{\rm{decoder}}}
\newcommand{\precons}{p_{\rm{reconstruct}}}

\newcommand{\laplace}{\mathrm{Laplace}} % Laplace distribution

\newcommand{\E}{\mathbb{E}}
\newcommand{\Ls}{\mathcal{L}}
\newcommand{\R}{\mathbb{R}}
\newcommand{\emp}{\tilde{p}}
\newcommand{\lr}{\alpha}
\newcommand{\reg}{\lambda}
\newcommand{\rect}{\mathrm{rectifier}}
\newcommand{\softmax}{\mathrm{softmax}}
\newcommand{\sigmoid}{\sigma}
\newcommand{\softplus}{\zeta}
\newcommand{\KL}{D_{\mathrm{KL}}}
\newcommand{\Var}{\mathrm{Var}}
\newcommand{\standarderror}{\mathrm{SE}}
\newcommand{\Cov}{\mathrm{Cov}}
% Wolfram Mathworld says $L^2$ is for function spaces and $\ell^2$ is for vectors
% But then they seem to use $L^2$ for vectors throughout the site, and so does
% wikipedia.
\newcommand{\normlzero}{L^0}
\newcommand{\normlone}{L^1}
\newcommand{\normltwo}{L^2}
\newcommand{\normlp}{L^p}
\newcommand{\normmax}{L^\infty}

\newcommand{\parents}{Pa} % See usage in notation.tex. Chosen to match Daphne's book.

\DeclareMathOperator*{\argmax}{arg\,max}
\DeclareMathOperator*{\argmin}{arg\,min}

\DeclareMathOperator{\sign}{sign}
\DeclareMathOperator{\Tr}{Tr}
\let\ab\allowbreak


\usepackage{hyperref}
\usepackage{url}
\usepackage{graphicx}
\usepackage{booktabs}
\usepackage{amsmath,amssymb}
\usepackage{amsthm}
\usepackage{algorithm}
\usepackage{algpseudocode}
\usepackage{capt-of}
\usepackage{xcolor}

\newtheorem{definition}{Definition}
\newtheorem{lemma}{Lemma}
\newtheorem{proposition}{Proposition}
\newtheorem{corollary}{Corollary}

\newcommand{\enc}{\operatorname{encode}}
\newcommand{\dec}{\operatorname{decode}}
\newcommand{\cover}{\operatorname{cover}}
\newcommand{\tok}[1]{\texttt{#1}}

%%%%%%%%% TITLE - PLEASE UPDATE TITLE HERE
\title{Exact Byte-Level Probabilities from Tokenized \\ Language Models for Fill-in-the-Middle \\ and Cross-Vocabulary Ensembling}

\author{Ambitious AI Researcher\\
AI Research Institute\\
123 AI Avenue, ML City\\
\texttt{researcher@institute.ai}
}

\newcommand{\fix}{\marginpar{FIX}}
\newcommand{\new}{\marginpar{NEW}}

\begin{document}

\maketitle

\begin{abstract}
%%%%%%%%% ABSTRACT %%%%%%%%%
Tokenized language models (LMs) are trained on token sequences, yet users interact with them through raw text. We show that this mismatch induces \emph{tokenization bias}: conditioning on the canonical tokenization of a prompt implicitly conditions on the event that a token boundary falls at the end of the prompt, which can assign zero probability to the correct continuation even for a perfectly trained model. The failure is acute when prompts end mid-token, as in fill-in-the-middle (FIM) code completion, and the incompatible output spaces of different tokenizers also prevent the direct ensembling of LMs. We formalize valid and \emph{cover} encodings and prove a Byte-Token Representation Lemma that expresses the probability of any byte string exactly in terms of a tokenized LM. Building on a structural result showing that every cover encoding is determined by its final token, we derive an inference-time algorithm that converts any tokenized LM into a statistically equivalent byte-level LM without retraining, requiring one model call per generated byte after a short initialization. On synthetic Markov chains, the method exactly recovers transition probabilities that a tokenized model gets wrong. On the HumanEval random-span FIM benchmark with CodeLlama-7b, byte-level prediction raises greedy pass rates under the suffix-prefix-middle template from 45.0\% to 63.9\%, outperforming token healing and token alignment. Mapping models with different vocabularies into a shared byte space further enables simple probability averaging, improving over the best member model on MMLU, PIQA, NQ, TriviaQA and HumanEval, while degrading on GSM8K, where members differ widely in capability.
%%%%%%%%% ABSTRACT %%%%%%%%%
\end{abstract}

\section{Introduction}
\label{sec:intro}
%%%%%%%%% INTRODUCTION %%%%%%%%%
Virtually all state-of-the-art language models operate on tokens rather than on raw bytes or characters~\citep{radford2019gpt2,brown2020gpt3,touvron2023llama2,jiang2023mistral}. Subword tokenizers such as byte-pair encoding (BPE)~\citep{gage1994bpe,sennrich2016bpe,wang2020bbpe}, WordPiece~\citep{schuster2012wordpiece} and SentencePiece~\citep{kudo2018sentencepiece} compress text into shorter sequences, which reduces the quadratic cost of attention~\citep{vaswani2017attention} and lets a fixed context window cover more text. However, the deterministic map from bytes to tokens is not innocuous. A tokenized LM learns a distribution over the \emph{canonical} token sequences produced by its tokenizer, whereas the quantity of interest to a user is usually the distribution over the next \emph{characters} given a textual prompt.

These two distributions disagree in a systematic way that we call \emph{tokenization bias}. Consider the code completion example in Figure~\ref{fig:teaser} (left): the prompt ends with \tok{return 5 + a}, and the intended continuation is the variable \tok{aa}. Because the tokenizer would have merged \tok{aa} into a single token had the full word been present, the training data essentially never contains the token \tok{a} followed by the token \tok{a}. Conditioned on the tokenized prompt, the model therefore assigns (near) zero probability to the correct completion. Crucially, this is \emph{not} a modeling error: as we show in Section~\ref{sec:method}, even a model that perfectly fits the training distribution exhibits the bias, because the tokenized prompt implicitly asserts that a token boundary lies at the end of the prompt. Prompts that end mid-token are ubiquitous in fill-in-the-middle (FIM) code completion~\citep{bavarian2022fim,fried2023incoder,roziere2023codellama}, where the span to be infilled can start anywhere. Practitioners resort to heuristics such as token healing~\citep{lundberg2023tokenhealing,dagan2024tokenizer} and token alignment~\citep{athiwaratkun2024alignment}, which back up one or more prompt tokens and constrain generation, but do not compute the correct conditional distribution.

Tokenization also prevents a second, seemingly unrelated capability: \emph{ensembling} LMs with different vocabularies. The next-token distributions of Llama~2~\citep{touvron2023llama2}, Mistral~\citep{jiang2023mistral} and Yi~\citep{young2024yi} live on disjoint supports, so they cannot be averaged directly; existing approaches rely on output reranking~\citep{jiang2023llmblender}, vocabulary alignment~\citep{xu2024eva,huang2024deepen,yu2024gac} or additional training~\citep{wan2024fusellm}.

We address both problems with one tool: an exact, training-free conversion of any tokenized LM into a byte-level LM. Our key observation is that the probability of a byte string can be written as a sum over the probabilities of a small set of token sequences, which we call \emph{cover encodings}, provided that we restrict attention to \emph{valid} encodings, i.e., token sequences the tokenizer could actually produce. Because the byte-level distribution is shared across all tokenizers, models with different vocabularies become directly comparable, and any ensembling strategy can be applied (Figure~\ref{fig:teaser}, right).

\textbf{Contributions.}
\textbf{(i)} We formally characterize tokenization bias as conditioning on a token boundary, and illustrate on Markov chains that it arises even for an ideal model (Section~\ref{sec:bias}).
\textbf{(ii)} We prove a Byte-Token Representation Lemma that maps token-level probabilities to exact byte-level probabilities, and show that each cover encoding is determined by its last token (Section~\ref{sec:cover}). This yields an efficient backward search for cover encodings and a byte-level sampling algorithm that needs a single LM call per byte (Section~\ref{sec:sampling}).
\textbf{(iii)} On the HumanEval random-span FIM benchmark, byte-level prediction removes most of the degradation caused by mid-token prompts, improving greedy pass rates in SPM mode from 45.0\% to 63.9\% and outperforming token healing and token alignment; it also improves PSM mode (Section~\ref{sec:exp_fim}).
\textbf{(iv)} We ensemble open-weight LMs with distinct tokenizers by averaging their byte-level predictions, outperforming the best individual model on several knowledge, reasoning and coding benchmarks, and we document when ensembling fails (Section~\ref{sec:exp_ensemble}).

\begin{figure}[t]
\begin{center}
\includegraphics[width=\linewidth]{figures/figure_1.png}
\end{center}
\caption{Left: tokenized LMs suffer from tokenization bias when prompts end mid-token, as in this code completion example, where the correct continuation \tok{a} receives zero probability. Our method uses the same model but predicts the correct byte. Right: our method maps the next-token predictions of arbitrary tokenized LMs to statistically equivalent next-byte predictions (Section~\ref{sec:method}), which enables any ensembling strategy, such as averaging or mixtures of experts.}
\label{fig:teaser}
\end{figure}
%%%%%%%%% INTRODUCTION %%%%%%%%%

\section{Related Work}
\label{sec:related_work}
%%%%%%%%% RELATED WORK %%%%%%%%%
\textbf{Tokenization and its side effects.}
Subword tokenization~\citep{sennrich2016bpe,schuster2012wordpiece,kudo2018sentencepiece,wang2020bbpe} is the default for modern LMs; see \citet{mielke2021between} for a historical overview. A growing literature documents its side effects, including degraded multilingual performance~\citep{rust2021tokenizer}, unfair cost across languages~\citep{petrov2023unfairness}, sensitivity of arithmetic to digit tokenization~\citep{singh2024tokenization}, and the influence of tokenizer choice on downstream quality~\citep{ali2024tokenizer,dagan2024tokenizer}. Theoretical work has analyzed BPE as an optimization problem~\citep{zouhar2023formal}, related tokenization quality to information-theoretic efficiency~\citep{zouhar2023noiseless}, studied the role of tokenization when learning Markov sources~\citep{rajaraman2024theory}, and formalized when token-level models consistently represent distributions over strings~\citep{gastaldi2024foundations}. \citet{phan2024understanding} study tokenization bias in Markov chains and propose a correction procedure; our work focuses on an efficient exact conversion that applies to arbitrary byte strings, FIM prompts, and cross-vocabulary ensembles of large models.

\textbf{Marginalizing over tokenizations.}
Since a string admits many token sequences, several works marginalize over tokenizations when \emph{scoring} text~\citep{cao2021marginal,chirkova2023marginalize}, typically with importance sampling, and subword regularization~\citep{kudo2018subword,provilkov2020bpedropout} trains models on non-canonical segmentations. Our perspective differs in two respects. First, for deterministic tokenizers, non-canonical (invalid) encodings have zero probability under the data distribution, so we restrict the sum to valid encodings, which makes it exact and tractable. Second, we target \emph{next-byte prediction} and generation, rather than likelihood evaluation of complete strings.

\textbf{Prompt boundary heuristics.}
Token healing~\citep{lundberg2023tokenhealing,dagan2024tokenizer} removes the last prompt token and constrains the first generated token to start with the removed characters. Token alignment~\citep{athiwaratkun2024alignment} backs up several tokens and constrains generation to match the removed characters. Both mitigate mid-token prompts, but neither accounts for all valid encodings nor weights them by their probabilities. Our method can be viewed as the principled limit of these heuristics: it marginalizes over every token boundary position that is consistent with the prompt.

\textbf{Byte-level and token-free models.}
Character- and byte-level LMs avoid tokenization altogether~\citep{alrfou2019character,xue2022byt5,clark2022canine,tay2022charformer,yu2023megabyte,wang2024mambabyte}, at the cost of longer sequences and the need to train new models from scratch. Our approach is complementary: it obtains byte-level predictions from existing tokenized models at inference time, without retraining.

\textbf{Ensembling and fusing LMs.}
Ensembles are a classic way to improve accuracy and calibration~\citep{dietterich2000ensemble,lakshminarayanan2017deepensembles}. For LMs, ensembling is hindered by vocabulary mismatch. Existing approaches select or fuse complete outputs~\citep{jiang2023llmblender,wang2023selfconsistency}, align vocabularies through learned or heuristic mappings~\citep{xu2024eva,huang2024deepen,yu2024gac}, optimize mixture weights at test time~\citep{mavromatis2024packllm}, distill multiple models into one~\citep{wan2024fusellm}, or transfer a model to a new tokenizer~\citep{minixhofer2024zeroshot}. Mixture-of-experts architectures~\citep{shazeer2017moe,jiang2024mixtral} share a single vocabulary by construction. In contrast, our byte-level conversion provides an exact common output space, so any combination rule can be applied at every generation step without training or heuristic alignment.
%%%%%%%%% RELATED WORK %%%%%%%%%

\section{Methodology}
\label{sec:method}
%%%%%%%%% METHOD %%%%%%%%%
\subsection{Setup and notation}
\label{sec:setup}
Let $\gA$ be the alphabet of bytes and $\gA^*$ the set of finite byte strings. We write $x_i^j=(x_i,\dots,x_j)$ for a substring and $|s|$ for the length of a string $s$. A tokenizer has a vocabulary $\gV\subset\gA^*$ of tokens, each of which is a non-empty byte string of length at most $L$, a deterministic encoder $\enc:\gA^*\to\gV^*$, and a decoder $\dec:\gV^*\to\gA^*$ that concatenates the bytes of the tokens. We denote token sequences by $\vt=(t_1,\dots,t_k)=t_1^k$ and concatenation by $\oplus$. We consider tokenizers that are lossless, $\dec(\enc(s))=s$, and focus on BPE~\citep{sennrich2016bpe} and maximum prefix encoding (MPE, i.e., greedy longest-match as in WordPiece~\citep{schuster2012wordpiece}).

Documents are byte strings drawn from an unknown distribution, and we write $P(x_1^n)$ for the probability that a document begins with $x_1^n$, so that $P(x_{n+1}\mid x_1^n)=P(x_1^{n+1})/P(x_1^n)$. A tokenized LM is trained on $\enc(\text{document})$ and defines next-token probabilities $P(t_{k+1}\mid t_1^k)$. We write $P(\vt)=\prod_{i=1}^{k}P(t_i\mid t_1^{i-1})$ for the probability that the token sequence begins with $\vt$. For the theory, we consider an \emph{ideal} tokenized LM whose distribution equals the distribution of $\enc(\text{document})$; in practice, we apply the same formulas to a trained model.

\subsection{Tokenization bias}
\label{sec:bias}
The standard way to obtain next-character predictions from a tokenized LM is to tokenize the prompt and read off the distribution of the first byte of the next token:
\begin{equation}
\label{eq:naive}
\hat P(x_{n+1}=a\mid x_1^n) \;=\; \sum_{\tau\in\gV:\,\tau_1=a} P\big(\tau \mid \enc(x_1^n)\big).
\end{equation}
Let $B_n$ denote the event that the canonical tokenization of the document has a token boundary right after position $n$. For the tokenizers we consider (see Proposition~\ref{prop:prefix} below), a document that begins with $x_1^n$ and has a boundary at $n$ is tokenized as $\enc(x_1^n)$ followed by further tokens. Hence, even for an ideal model, Equation~\ref{eq:naive} computes
\begin{equation}
\label{eq:bias}
\hat P(x_{n+1}=a\mid x_1^n) \;=\; P\big(x_{n+1}=a \mid x_1^n,\, B_n\big)\;\neq\; P\big(x_{n+1}=a\mid x_1^n\big)\quad\text{in general.}
\end{equation}
We call the discrepancy in Equation~\ref{eq:bias} \emph{tokenization bias}. It arises whenever knowing that a token ends at position $n$ is informative about the next byte, which is precisely the case when the prompt ends with a token that is a prefix of other tokens.

\begin{figure}[t]
\begin{center}
\includegraphics[width=\linewidth]{figures/figure_2.png}
\end{center}
\caption{Tokenization bias on a 1st-order Markov chain over $\{\tok{A},\tok{B}\}$ with vocabulary $\{\tok{A},\tok{B},\tok{AA}\}$. Given the context token ``\tok{A}'', the tokenized model never samples ``\tok{A}'' as the next byte, whereas the true probability is $1-\alpha$. In practice, the same bias occurs when prompts end with tokens that are prefixes of other tokens, which is common in FIM tasks.}
\label{fig:markov1}
\end{figure}

\textbf{Example.} Figure~\ref{fig:markov1} shows a 1st-order Markov chain over $\{\tok{A},\tok{B}\}$ that stays at \tok{A} with probability $1-\alpha$ and at \tok{B} with probability $1-\beta$. With vocabulary $\{\tok{A},\tok{B},\tok{AA}\}$, both BPE and MPE merge every pair of consecutive \tok{A}s, so the string \tok{AABABAA} is tokenized as \tok{AA|B|A|B|AA}. A standalone token \tok{A} can therefore only be followed by \tok{B}: the induced token process assigns $P(\tok{A}\mid\tok{A})=P(\tok{AA}\mid\tok{A})=0$ and $P(\tok{B}\mid\tok{A})=1$. A prompt ending in \tok{BA} is encoded with a final token \tok{A}, and Equation~\ref{eq:naive} yields $\hat P(\tok{A}\mid \dots\tok{BA})=0$, although the true probability is $1-\alpha$. No amount of training removes this error, because it is a property of the token process itself.

\subsection{Valid encodings, cover encodings, and the Byte-Token Representation Lemma}
\label{sec:cover}
Our approach replaces conditioning on one tokenization with an exact marginalization. The key is to know which token sequences can occur at all.

\begin{definition}[Valid encoding]
\label{def:valid}
A token sequence $\vt\in\gV^*$ is \emph{valid} if $\enc(\dec(\vt))=\vt$, and \emph{invalid} otherwise.
\end{definition}

Under the ideal model, invalid sequences are never generated, $P(\vt)=0$, because they never appear in tokenized training data. In Figure~\ref{fig:markov1}, \tok{A|A} is invalid since $\enc(\tok{AA})=\tok{AA}$. We rely on the following property, proven in Appendix~\ref{app:proofs}.

\begin{proposition}[Prefix closure]
\label{prop:prefix}
For BPE and MPE, every prefix of a valid encoding is valid: if $t_1^k$ is valid, then $t_1^j$ is valid and equals $\enc(\dec(t_1^j))$ for all $j\leq k$.
\end{proposition}

\begin{definition}[Cover encodings]
\label{def:cover}
The \emph{cover encodings} of a byte string $x_1^n$ are the valid token sequences $\vt=t_1^k$ such that $x_1^n$ is a prefix of $\dec(\vt)$ and $\dec(t_1^{k-1})$ is a strict prefix of $x_1^n$. We denote this set by $\cover(x_1^n)$.
\end{definition}

In words, a cover encoding tokenizes the string up to some position $j<n$ and then ends with a single token that covers the remaining bytes $x_{j+1}^n$ and may extend beyond position $n$ (Figure~\ref{fig:cover}, left). The canonical encoding $\enc(x_1^n)$ is always a cover encoding.

\begin{lemma}[Byte-Token Representation]
\label{lem:btr}
Under the ideal tokenized LM, for every byte string $x_1^n$,
\begin{equation}
\label{eq:btr}
P(x_1^n) \;=\; \sum_{\vt\in\cover(x_1^n)} P(\vt).
\end{equation}
\end{lemma}
\begin{proof}[Proof sketch]
Let $T=\enc(\text{document})$. If the document begins with $x_1^n$, let $k^*$ be the smallest index such that $|\dec(T_1^{k^*})|\geq n$. By Proposition~\ref{prop:prefix}, $T_1^{k^*}$ is valid, and by construction it belongs to $\cover(x_1^n)$. Conversely, if $T$ begins with some $\vt\in\cover(x_1^n)$, the document begins with $x_1^n$. Finally, no two distinct cover encodings can both be prefixes of $T$, since the shorter one would already reach length $n$, contradicting the definition for the longer one. Hence the event that the document begins with $x_1^n$ is the disjoint union of the events that $T$ begins with $\vt$, $\vt\in\cover(x_1^n)$. A full proof is in Appendix~\ref{app:proofs}.
\end{proof}

Lemma~\ref{lem:btr} turns any tokenized LM into a byte-level LM: next-byte probabilities follow from $P(x_{n+1}\mid x_1^n)=P(x_1^{n+1})/P(x_1^n)$, and each $P(\vt)$ is computed by the chain rule with the LM. It remains to enumerate the cover encodings efficiently.

\begin{figure}[t]
\begin{center}
\includegraphics[width=\linewidth]{figures/figure_4.png}
\end{center}
\caption{Left: representation of $P(x_1^n)$ with a tokenized LM, with examples of cover encodings and of valid and invalid encodings. Right: illustration of Algorithm~\ref{alg:cover} on the query string \tok{yummy} with a BPE vocabulary. Green ticks and red crosses denote valid and invalid encodings (Definition~\ref{def:valid}). The search terminates once the suffix exceeds the maximum token length $L$.}
\label{fig:cover}
\end{figure}

\textbf{Structure of cover encodings.} Naively, the number of token sequences that decode to a string grows exponentially with its length. Validity collapses this set dramatically:

\begin{proposition}[Cover encodings are determined by their last token]
\label{prop:structure}
Every $\vt\in\cover(x_1^n)$ has the form $\vt=\enc(x_1^{j})\oplus\tau$, where $n-L\leq j<n$ and $\tau\in\gV$ has $x_{j+1}^n$ as a prefix. Consequently, $|\cover(x_1^n)|\leq\sum_{i=0}^{L-1}\big|\{\tau\in\gV:\ x_{n-i}^n \text{ is a prefix of } \tau\}\big|$.
\end{proposition}
\begin{proof}
Write $\vt=t_1^k$ and $\dec(t_1^{k-1})=x_1^j$ with $j<n$. By Proposition~\ref{prop:prefix}, $t_1^{k-1}$ is valid, so $t_1^{k-1}=\enc(x_1^j)$. The last token $t_k$ must contain $x_{j+1}^n$ as a prefix, so $n-j\leq|t_k|\leq L$.
\end{proof}

Proposition~\ref{prop:structure} implies that we never need to search over the full history: it suffices to look backwards at most $L$ bytes, and for each candidate boundary $j$ the prefix is fixed to the canonical encoding $\enc(x_1^j)$. This gives Algorithm~\ref{alg:cover}, illustrated in Figure~\ref{fig:cover} (right). The probability of each cover encoding factorizes as $P(\enc(x_1^j)\oplus\tau)=P(\enc(x_1^j))\,P(\tau\mid\enc(x_1^j))$, so all of them can be scored with at most $L$ forward passes, one per candidate boundary, which can be batched and share a key-value cache whenever the canonical prefixes coincide.

\begin{algorithm}[t]
\caption{Cover encoding search by backward tokenization}
\label{alg:cover}
\begin{algorithmic}[1]
\Require byte string $x_1^n$, vocabulary $\gV$ with maximum token length $L$, tokenizer $\enc$, tokenized LM $P$
\State $\gC\gets\emptyset$
\For{$i=0,1,\dots,\min(n,L)-1$} \Comment{the last token starts at position $n-i$}
  \State $\vp\gets\enc(x_1^{n-i-1})$ \Comment{canonical encoding of the remaining prefix}
  \For{each $\tau\in\gV$ such that $x_{n-i}^n$ is a prefix of $\tau$}
    \If{$\enc(\dec(\vp\oplus\tau))=\vp\oplus\tau$} \Comment{validity check (Definition~\ref{def:valid})}
      \State $\gC\gets\gC\cup\{\vp\oplus\tau\}$ with score $P(\vp)\,P(\tau\mid\vp)$
    \EndIf
  \EndFor
\EndFor
\State \Return $\gC=\cover(x_1^n)$ and the scores $\{P(\vt)\}_{\vt\in\gC}$
\end{algorithmic}
\end{algorithm}

\subsection{Exact next-byte prediction and sampling}
\label{sec:sampling}
Algorithm~\ref{alg:cover} is only needed once, for the prompt. During generation, the cover encodings of $x_1^{n+1}$ can be obtained from those of $x_1^n$ with a single LM call. Partition $\cover(x_1^n)$ into the encodings that extend strictly beyond position $n$ and the unique encoding that ends exactly at $n$, which by validity is $\vt^\star=\enc(x_1^n)$. For each byte $a\in\gA$, define
\begin{align}
\gC_{n+1}(a) &= \big\{\vt\in\cover(x_1^n):\ |\dec(\vt)|>n,\ \dec(\vt)_{n+1}=a\big\}, \label{eq:regroup}\\
\bar\gC_{n+1}(a) &= \big\{\vt^\star\oplus\tau:\ \tau\in\gV,\ \tau_1=a,\ \vt^\star\oplus\tau\text{ is valid}\big\}. \label{eq:extend}
\end{align}

\begin{proposition}[Recursive cover update]
\label{prop:recursion}
For every $a\in\gA$, $\cover(x_1^n a)=\gC_{n+1}(a)\,\cup\,\bar\gC_{n+1}(a)$, the union being disjoint. Therefore, under the ideal model,
\begin{equation}
\label{eq:nextbyte}
P(x_{n+1}=a\mid x_1^n)\;=\;\frac{\sum_{\vt\in\gC_{n+1}(a)}P(\vt)\;+\;P(\vt^\star)\sum_{\vt^\star\oplus\tau\in\bar\gC_{n+1}(a)}P(\tau\mid\vt^\star)}{\sum_{b\in\gA}\Big[\sum_{\vt\in\gC_{n+1}(b)}P(\vt)\;+\;P(\vt^\star)\sum_{\vt^\star\oplus\tau\in\bar\gC_{n+1}(b)}P(\tau\mid\vt^\star)\Big]}.
\end{equation}
\end{proposition}

The first term in the numerator reuses probabilities that were already computed for $\cover(x_1^n)$; only the second term requires a new forward pass, on $\vt^\star$. Under the ideal model, the denominator equals $P(x_1^n)$ by Lemma~\ref{lem:btr}; for a trained model, which places a small amount of mass on invalid continuations, normalizing over $b$ discards that mass consistently. Comparing Equation~\ref{eq:nextbyte} with Equation~\ref{eq:naive} makes the source of tokenization bias explicit: the naive estimate keeps only the second term and ignores all encodings in which a token straddles position $n$.

The full procedure (Algorithm~\ref{alg:sampling} in Appendix~\ref{app:algorithms}, illustrated in Figure~\ref{fig:sampling}) proceeds in five steps: (1) obtain $\cover(x_1^n)$, from Algorithm~\ref{alg:cover} at the first step or from the previous iteration otherwise; (2) regroup the encodings that extend beyond $n$ by their $(n{+}1)$-th byte (Equation~\ref{eq:regroup}); (3) run the LM once on $\vt^\star=\enc(x_1^n)$, accept or reject each next token through the validity check, and group the accepted ones by their first byte (Equation~\ref{eq:extend}); (4) combine both groups into $\cover(x_1^n a)$ for every byte $a$ and compute Equation~\ref{eq:nextbyte}; (5) sample (or select greedily) $x_{n+1}$ and keep $\cover(x_1^{n+1})$ for the next step. Note that $\enc(x_1^{n+1})$ is itself a member of $\cover(x_1^{n+1})$, so its probability is available at the next iteration. All computations are carried out in log space.

\begin{figure}[t]
\begin{center}
\includegraphics[width=\linewidth]{figures/figure_5.png}
\end{center}
\caption{Illustration of sampling $x_{n+1}$ from a tokenized LM $P(t_{k+1}\mid t_1^k)$, continuing the example of Figure~\ref{fig:cover}. Cover encodings longer than $n$ are regrouped by their next byte (step 2), and the canonical encoding is extended by one LM call with accept/reject validity checks (step 3); the two groups are combined (step 4) before sampling (step 5). Probabilities $P(\vt)$ are omitted for simplicity.}
\label{fig:sampling}
\end{figure}

\textbf{Cost.} After initialization (at most $L$ batched forward passes), the method requires one LM call per generated \emph{byte}, compared with one call per \emph{token} for standard decoding; the overhead thus scales with the average number of bytes per token. The validity checks only require the tokenizer, not the LM, and only tokens with non-negligible probability affect Equation~\ref{eq:nextbyte}.

\subsection{Applications}
\label{sec:applications}
\textbf{Fill-in-the-middle.} FIM models are prompted with sentinel tokens, either in prefix-suffix-middle (PSM) order, \texttt{<PRE>}\,prefix\,\texttt{<SUF>}\,suffix\,\texttt{<MID>}, or in suffix-prefix-middle (SPM) order, in which the prompt ends with the prefix and the middle continues it directly~\citep{bavarian2022fim,roziere2023codellama}. In SPM, the prompt therefore ends mid-token whenever the span to be infilled starts inside a word. We treat sentinel tokens as atomic boundaries, so the cover search never crosses them, and apply Algorithm~\ref{alg:cover} to the text that follows the last sentinel. End-of-sequence tokens are treated as an additional symbol of the byte alphabet, so generation stops when this symbol is selected.

\textbf{Byte-level greedy and beam search.} Equation~\ref{eq:nextbyte} defines a proper autoregressive distribution over bytes, so standard decoding strategies apply. Greedy decoding selects $\arg\max_a P(x_{n+1}{=}a\mid x_1^n)$, which aggregates the mass of \emph{all} tokens starting with byte $a$ and is therefore a different decision rule from token-level greedy decoding. For beam search with width $B$, each hypothesis is a byte string with its own cover set, scored by $\log P(x_1^n)$.

\textbf{Cross-vocabulary ensembling.} Because Equation~\ref{eq:nextbyte} lives in the common space $\gA$, predictions of models $P_1,\dots,P_M$ with different tokenizers can be combined directly. We use the simplest rule, the average $P_{\text{ens}}(x_{n+1}\mid x_1^n)=\frac{1}{M}\sum_{m=1}^M P_m(x_{n+1}\mid x_1^n)$. After a byte is selected, every member updates its own cover set with Proposition~\ref{prop:recursion}. Other rules, such as weighted averaging or mixtures of experts, can be plugged in without modification.
%%%%%%%%% METHOD %%%%%%%%%

\section{Experiments}
\label{sec:experiments}
%%%%%%%%% EXPERIMENTS %%%%%%%%%
We evaluate (i) whether the byte-token conversion recovers the true byte-level distribution in a controlled setting, (ii) whether it fixes the mid-token failure mode in FIM code completion, and (iii) whether the byte space enables effective ensembling of models with different vocabularies.

\subsection{Setup}
\label{sec:exp_setup}
\textbf{Models.} For the synthetic study, we use 1st- and 3rd-order Markov chains over $\{\tok{A},\tok{B}\}$ and a GPT-2 architecture~\citep{radford2019gpt2} with 6 hidden layers trained on the tokenized output of the 3rd-order chain. For FIM, we use CodeLlama-7b~\citep{roziere2023codellama}. For general-purpose evaluation, we use Llama2-7B~\citep{touvron2023llama2}, Yi-1.5-6B~\citep{young2024yi} and Mistral-7B-v0.3~\citep{jiang2023mistral}, and for code generation CodeLlama2-7b, CodeLlama2-13b~\citep{roziere2023codellama} and Yi-Coder-1.5B. All are Transformer-based LMs and are used without any weight updates.

\textbf{Tasks and metrics.} For FIM, we use the HumanEval random-span infilling benchmark~\citep{bavarian2022fim} in both PSM and SPM modes, reporting greedy pass rates and pass@10, where pass@10 is the maximum over a temperature sweep in $[0,1]$. For knowledge and reasoning, we use MMLU~\citep{hendrycks2021mmlu}, PIQA~\citep{bisk2020piqa}, Natural Questions (NQ)~\citep{kwiatkowski2019nq}, TriviaQA~\citep{joshi2017triviaqa} and GSM8K~\citep{cobbe2021gsm8k}, reporting accuracy. For code generation, we use HumanEval~\citep{chen2021humaneval} and MBPP~\citep{austin2021mbpp}, reporting pass@1. For each model, we compare predictions from the standard token-level distribution (\emph{Token}) with predictions from our byte-level distribution (\emph{Byte}) obtained from the same model.

\textbf{Baselines.} Besides standard token-level prediction, we compare against token healing~\citep{lundberg2023tokenhealing,dagan2024tokenizer} and token alignment~\citep{athiwaratkun2024alignment} for FIM, both applied in SPM mode where they are relevant. For ensembling, we compare against the individual member models and, for the multiple-choice tasks MMLU and PIQA, against top-2 majority voting. All ensembles combine the top-2 individual models per benchmark by averaging their byte-level probabilities.

\subsection{Tokenization bias on synthetic data}
\label{sec:exp_synthetic}
\begin{figure}[t]
\begin{minipage}[c]{0.31\linewidth}
\begin{center}
\includegraphics[width=\linewidth]{figures/figure_3.png}
\end{center}
\vspace{-6pt}
\caption{Tokenization bias on a 3rd-order Markov chain. Our byte-token conversion accurately recovers $P(x_{n+1}\mid x_{n-2}^n)$ of the chain.}
\label{fig:markov3}
\end{minipage}
\hfill
\begin{minipage}[c]{0.65\linewidth}
\captionof{table}{Code generation (pass@1, \%). The ensemble averages the byte-level predictions of the top-2 models (Yi-Coder-1.5B and CodeLlama2-13b).}
\label{tab:code}
\begin{center}
\small
\setlength{\tabcolsep}{4pt}
\begin{tabular}{lcccc}
\toprule
 & \multicolumn{2}{c}{HumanEval} & \multicolumn{2}{c}{MBPP} \\
\cmidrule(lr){2-3}\cmidrule(lr){4-5}
Model & Token & Byte & Token & Byte \\
\midrule
CodeLlama2-7b & 32.3 & 28.7 & 40.6 & 42.4 \\
CodeLlama2-13b & 35.8 & 33.5 & 47.6 & 47.4 \\
Yi-Coder-1.5B & 38.4 & 36.5 & 52.8 & \textbf{53.6} \\
\midrule
Top-2 ensemble (ours) & -- & \textbf{42.1} & -- & \textbf{53.6} \\
\bottomrule
\end{tabular}
\end{center}
\vspace{2pt}
\captionof{table}{Beam width ablation for byte-level decoding on HumanEval random-span FIM with CodeLlama-7b (pass rate, \%).}
\label{tab:beam}
\begin{center}
\small
\begin{tabular}{lccc}
\toprule
Beam width & 1 & 2 & 4 \\
\midrule
PSM & 65.7 & 66.8 & \textbf{67.7} \\
SPM & 63.9 & 65.9 & \textbf{66.8} \\
\bottomrule
\end{tabular}
\end{center}
\end{minipage}
\end{figure}

In the 1st-order chain of Figure~\ref{fig:markov1}, the token-level model assigns probability $0.0$ to \tok{A} following a context that ends in the token \tok{A}, whereas the ground-truth probability is $1-\alpha$, confirming the analysis of Section~\ref{sec:bias}. Figure~\ref{fig:markov3} reports the 3rd-order chain, for which we trained the GPT-2 model on tokenized sequences. For each of the eight 3-byte contexts, we compare the ground-truth probability that the next byte is \tok{A}, the naive estimate $P(x_{n+1}\mid t_1^k)$ of Equation~\ref{eq:naive}, and our byte-level estimate. The naive estimate is severely biased: for most contexts it places (close to) zero probability on \tok{A} even though the ground truth is substantial (e.g., above $0.6$ for \tok{BAB}, \tok{ABB} and \tok{ABA}). In contrast, our bias-corrected estimate, computed from the very same trained model, closely matches the ground truth for all contexts. This confirms both that tokenization bias is severe for trained models and that Lemma~\ref{lem:btr} recovers the original transition probabilities without retraining.

\subsection{Fill-in-the-middle code completion}
\label{sec:exp_fim}
\begin{figure}[t]
\begin{center}
\includegraphics[width=\linewidth]{figures/figure_6.png}
\end{center}
\caption{Pass rates on the HumanEval random-span benchmark~\citep{bavarian2022fim} with CodeLlama-7b for infilling Python. Left: greedy pass rates and pass@10 (maximum over the temperature sweep on the right) for the PSM and SPM templates; token healing~\citep{dagan2024tokenizer} and alignment~\citep{athiwaratkun2024alignment} are applied in SPM mode. Right: pass@10 as a function of temperature. The SPM template forces predictions from the middle of a token, which byte-level prediction handles successfully.}
\label{fig:fim}
\end{figure}

Figure~\ref{fig:fim} summarizes the FIM results. \textbf{Mid-token prompts break token-level prediction.} Switching from PSM to SPM reduces the token-level greedy pass rate from 64.9\% to 45.0\% and pass@10 from 84.0\% to 66.5\%. The right panel shows that the gap persists across all sampling temperatures, so it cannot be fixed by sampling more diversely; the issue is that the correct continuation is out of distribution given the canonical prompt tokenization. \textbf{Byte-level prediction removes the degradation.} With byte-level predictions from the same model, the SPM greedy pass rate rises to 63.9\% (+18.9 points) and pass@10 to 84.3\% (+17.8 points), on par with PSM; the byte-level pass@10 curves for PSM and SPM nearly coincide across temperatures. \textbf{Comparison with heuristics.} Token healing (62.7\% greedy, 83.8\% pass@10) and token alignment (63.0\%, 82.9\%) recover most of the gap, confirming that the SPM failure stems from prompt boundaries, but both remain below byte-level prediction in SPM mode (63.9\%, 84.3\%). \textbf{Gains also in PSM.} Even in PSM mode, where prompts end with a sentinel token and tokenization issues should be minimal, byte-level greedy decoding improves over token-level greedy decoding (65.7\% vs.\ 64.9\%), while pass@10 is essentially unchanged (84.1\% vs.\ 84.0\%). A possible explanation is that byte-level greedy decoding aggregates the probability of all tokens sharing a first byte (Section~\ref{sec:applications}), and that the middle span itself may start inside a word of the prefix; we did not isolate these factors.

\textbf{Error analysis.} Inspecting failures of the token-level model in SPM mode, we find cases in which the prompt ends in the middle of an identifier, for example in variables such as \tok{delimeter}, \tok{collatz} and \tok{palindrome}. The token-level model fails to complete the word because the remaining characters do not form a token that can follow the canonical prompt tokenization, while the byte-level method identifies the correct completion.

\begin{table}[t]
\caption{Knowledge and reasoning benchmarks (accuracy, \%). For each model, Token and Byte denote predictions from the token-level distribution and from our byte-level conversion of the same model. Voting and the ensemble combine the top-2 individual models per benchmark; voting is only applicable to the multiple-choice tasks. Best result per benchmark in bold.}
\label{tab:general}
\begin{center}
\small
\setlength{\tabcolsep}{4.2pt}
\begin{tabular}{lcccccccccc}
\toprule
 & \multicolumn{2}{c}{MMLU} & \multicolumn{2}{c}{PIQA} & \multicolumn{2}{c}{NQ} & \multicolumn{2}{c}{TriviaQA} & \multicolumn{2}{c}{GSM8K} \\
\cmidrule(lr){2-3}\cmidrule(lr){4-5}\cmidrule(lr){6-7}\cmidrule(lr){8-9}\cmidrule(lr){10-11}
Model & Token & Byte & Token & Byte & Token & Byte & Token & Byte & Token & Byte \\
\midrule
Llama2-7B & 45.4 & 45.7 & 78.1 & 78.2 & 25.0 & 23.4 & 58.4 & 58.1 & 15.1 & 12.5 \\
Yi-1.5-6B & 63.4 & 63.4 & 78.5 & 78.5 & 22.8 & 22.7 & 53.7 & 53.4 & 61.3 & \textbf{61.5} \\
Mistral-7B-v0.3 & 62.1 & 62.1 & 80.1 & 80.3 & 28.5 & 28.8 & 63.6 & 63.5 & 39.0 & 39.0 \\
\midrule
Voting (top-2) & -- & 62.1 & -- & 80.1 & -- & -- & -- & -- & -- & -- \\
Top-2 ensemble (ours) & -- & \textbf{65.4} & -- & \textbf{80.7} & -- & \textbf{30.0} & -- & \textbf{64.2} & -- & 55.8 \\
\bottomrule
\end{tabular}
\end{center}
\end{table}

\textbf{Ablation: byte-level beam search.} Table~\ref{tab:beam} varies the beam width of byte-level decoding on the FIM task. Wider beams yield consistent gains in both modes: from 65.7\% to 67.7\% in PSM (+2.0 points) and from 63.9\% to 66.8\% in SPM (+2.9 points) when going from width 1 (greedy) to width 4. Hence the byte-level distribution is compatible with standard search-based decoding, and its benefits are complementary to those of beam search.

\subsection{Byte-level conversion and cross-vocabulary ensembling}
\label{sec:exp_ensemble}


Tables~\ref{tab:general} and~\ref{tab:code} report token-level and byte-level results for individual models, and the results of ensembling the top-2 models of each benchmark in byte space.

\textbf{Individual models.} On benchmarks whose prompts end at natural boundaries, there is no mid-token issue to fix, and byte-level and token-level results are close on most tasks: MMLU, PIQA and TriviaQA differ by at most 0.3 points for every model, and Mistral-7B-v0.3 is within 0.3 points on all five benchmarks. There are, however, notable exceptions in both directions. Byte-level prediction reduces accuracy for Llama2-7B on NQ ($-1.6$) and GSM8K ($-2.6$), and reduces HumanEval pass@1 for all three code models ($-3.6$, $-2.3$ and $-1.9$ points), whereas it improves MBPP for CodeLlama2-7b ($+1.8$) and Yi-Coder-1.5B ($+0.8$). For an ideal model, byte-level and token-level \emph{sampling} define the same distribution over strings; the differences therefore stem from decoding rules and model imperfections. In particular, greedy decoding over bytes is not equivalent to greedy decoding over tokens, and the byte-level conversion discards the mass that a trained model places on invalid encodings. We have not disentangled these effects, and the HumanEval regressions indicate that byte-level greedy decoding is not uniformly beneficial for open-ended generation.

\textbf{Ensembles across vocabularies.} The byte space allows us to average models whose tokenizers are incompatible. On MMLU, the ensemble of Yi-1.5-6B and Mistral-7B-v0.3 reaches 65.4\%, improving by 2.0 points over the best member (63.4\%) and by 3.3 points over top-2 voting (62.1\%). The ensemble also outperforms the best individual result on PIQA (80.7\% vs.\ 80.3\%), NQ (30.0\% vs.\ 28.8\%) and TriviaQA (64.2\% vs.\ 63.6\%). For code, the ensemble of Yi-Coder-1.5B and CodeLlama2-13b achieves 42.1\% on HumanEval, 3.7 points above the best individual model (38.4\%), even though each member's byte-level score is lower than its token-level score; on MBPP it matches the best member (53.6\%).

\textbf{Limitation of ensembling.} Ensembling does not help when members differ widely in capability. On GSM8K, Yi-1.5-6B outperforms Mistral-7B-v0.3 by more than 20 points (61.5\% vs.\ 39.0\% in byte space), and the uniform average (55.8\%) falls 5.7 points below the stronger model. In multi-step generation such as math reasoning, a single poor byte choice induced by the weaker model can derail the solution, so uniform averaging is ill-suited to unequal members; weighted or input-dependent combination rules are natural remedies that our framework supports but that we have not evaluated.

\textbf{Statistical caveats.} All numbers are single evaluations on the standard test sets; we do not report confidence intervals. Several differences, especially those below one point (e.g., PIQA, TriviaQA, and the PSM pass@10), are small relative to the sampling variability expected on benchmarks of this size and should be interpreted with caution. The large effect in SPM mode (+18.9 greedy points) is unlikely to be explained by noise, whereas the moderate ensemble gains on MMLU (+2.0) and HumanEval (+3.7, i.e., about six of 164 problems) would benefit from confirmation with significance tests.
%%%%%%%%% EXPERIMENTS %%%%%%%%%

\section{Conclusion}
\label{sec:conclusion}
%%%%%%%%% CONCLUSION %%%%%%%%%
We identified tokenization bias as an intrinsic property of tokenized LMs: conditioning on the canonical tokenization of a prompt implicitly conditions on a token boundary at its end, which can make the correct continuation impossible. By restricting attention to valid encodings, we derived the Byte-Token Representation Lemma, showed that cover encodings are determined by their final token, and obtained an exact, training-free algorithm that turns any tokenized LM into a statistically equivalent byte-level LM with one model call per byte. Empirically, the conversion recovers ground-truth transition probabilities on Markov chains, removes most of the degradation caused by mid-token FIM prompts (45.0\%$\to$63.9\% greedy pass rate in SPM mode), outperforms token healing and token alignment, and enables simple averaging of models with different vocabularies, which improves over the best member on several benchmarks.

\textbf{Limitations and future work.} First, byte-level decoding requires one forward pass per byte rather than per token, which increases latency roughly by the average number of bytes per token; speculative or multi-byte decoding could reduce this cost. Second, our guarantees hold for an ideal model and for tokenizers satisfying prefix closure; production tokenizers with pre-tokenization rules or normalization require care, and trained models place non-zero mass on invalid encodings. Third, byte-level greedy decoding does not uniformly improve individual models, as the HumanEval results show, and uniform ensembling hurts when members are unequal, as on GSM8K. Finally, our evaluation reports single runs without confidence intervals. Promising directions include learned or adaptive ensemble weights, byte-level mixtures of experts over specialized open-weight models, and applying the conversion to other settings where prompt boundaries matter, such as constrained decoding.
%%%%%%%%% CONCLUSION %%%%%%%%%

\section*{Reproducibility Statement}
All quantities used by our method are defined in Section~\ref{sec:method}: Algorithm~\ref{alg:cover} specifies the cover search, Equation~\ref{eq:nextbyte} and Algorithm~\ref{alg:sampling} (Appendix~\ref{app:algorithms}) specify next-byte prediction, and Appendix~\ref{app:proofs} contains the proofs. The method requires no training and only uses publicly available model weights, tokenizers and benchmarks.

%%%%%%%%% REFERENCES %%%%%%%%%
\bibliography{references}
\bibliographystyle{iclr2025_conference}
%%%%%%%%% REFERENCES %%%%%%%%%

%%%%%%%%% Appendix %%%%%%%%%
\appendix
\section{Proofs}
\label{app:proofs}

Throughout, documents are finite byte strings, $T=\enc(\text{document})$ is the canonical token sequence of a random document, and $P(\vt)$ denotes the probability that $T$ begins with $\vt$, which is what an ideal tokenized LM computes by the chain rule.

\textbf{Proof of Proposition~\ref{prop:prefix} (prefix closure).}
Let $\vt=t_1^k=\enc(s)$ and fix $j\leq k$; write $u=\dec(t_1^j)$, so that $s=uv$ for some $v$. We show $\enc(u)=t_1^j$.

\emph{MPE.} MPE repeatedly emits the longest token that is a prefix of the remaining string. We argue by induction over the emitted tokens. Suppose the first $i<j$ tokens of $\enc(u)$ equal $t_1^i$, and let $r$ be the remaining part of $u$, so that $r v'$ is the remaining part of $s$ for some $v'$. On $s$, MPE emitted $t_{i+1}$, the longest token that is a prefix of $rv'$; since $t_{i+1}$ is also a prefix of $r$ (because $i+1\leq j$), and any token that is a prefix of $r$ is also a prefix of $rv'$ and hence no longer than $t_{i+1}$, MPE emits $t_{i+1}$ on $u$ as well. After $j$ steps $u$ is exhausted, so $\enc(u)=t_1^j$.

\emph{BPE.} BPE starts from single bytes and repeatedly merges the adjacent pair with the lowest merge rank, breaking ties from left to right, until no mergeable pair remains. Since merges only coarsen the segmentation and the final segmentation of $s$ has a boundary between $u$ and $v$, no merge in the trajectory on $s$ crosses that boundary. We claim that running BPE on $u$ performs exactly the subsequence of merges of the trajectory on $s$ that lie within $u$. Consider any state in which the parts within $u$ coincide. In the trajectory on $s$, merges within $v$ do not change the part within $u$; when the next merge within $u$ occurs, it is the lowest-ranked pair among all pairs of $s$, hence also among the pairs of $u$, and in case of ties it is the leftmost one, since all pairs of $u$ lie to the left of pairs involving $v$. This is exactly the merge BPE selects on $u$. Finally, when BPE terminates on $s$, no mergeable pair remains in $u$, so BPE on $u$ terminates at the same segmentation $t_1^j$. In both cases, $t_1^j$ is valid since $\enc(\dec(t_1^j))=\enc(u)=t_1^j$. \hfill$\square$

\textbf{Proof of Lemma~\ref{lem:btr}.}
Let $E$ be the event that the document begins with $x_1^n$ and, for $\vt\in\cover(x_1^n)$, let $E_{\vt}$ be the event that $T$ begins with $\vt$.
($E\subseteq\bigcup_{\vt}E_{\vt}$.) Suppose the document begins with $x_1^n$, and let $k^*$ be the smallest index with $|\dec(T_1^{k^*})|\geq n$, which exists because $\dec(T)$ is the document. Then $x_1^n$ is a prefix of $\dec(T_1^{k^*})$, $\dec(T_1^{k^*-1})$ is a strict prefix of $x_1^n$ by minimality, and $T_1^{k^*}$ is valid by Proposition~\ref{prop:prefix}. Hence $T_1^{k^*}\in\cover(x_1^n)$.
($\bigcup_{\vt}E_{\vt}\subseteq E$.) If $T$ begins with $\vt\in\cover(x_1^n)$, the document begins with $\dec(\vt)$, which has $x_1^n$ as a prefix.
(Disjointness.) If $T$ begins with two distinct $\vt,\vt'\in\cover(x_1^n)$, one of them, say $\vt$, is a strict prefix of the other, $\vt=t_1'^{\,k}$ with $k<|\vt'|$. Then $\dec(t_1'^{\,|\vt'|-1})$ has $\dec(\vt)$, and hence $x_1^n$, as a prefix, contradicting the requirement that it be a strict prefix of $x_1^n$.
Therefore $P(x_1^n)=P(E)=\sum_{\vt\in\cover(x_1^n)}P(E_{\vt})=\sum_{\vt\in\cover(x_1^n)}P(\vt)$. \hfill$\square$

\textbf{Proof of Proposition~\ref{prop:recursion}.}
Let $\vt=t_1^k\in\cover(x_1^na)$ and let $j=|\dec(t_1^{k-1})|\leq n$. If $j<n$, then $\dec(t_1^{k-1})$ is a strict prefix of $x_1^n$ and $x_1^n$ is a prefix of $\dec(\vt)$, so $\vt\in\cover(x_1^n)$; moreover $|\dec(\vt)|\geq n+1$ and $\dec(\vt)_{n+1}=a$, hence $\vt\in\gC_{n+1}(a)$. If $j=n$, then $t_1^{k-1}$ is valid and decodes to $x_1^n$, so $t_1^{k-1}=\enc(x_1^n)=\vt^\star$, and $t_k$ starts with $a$; hence $\vt\in\bar\gC_{n+1}(a)$. Conversely, every element of $\gC_{n+1}(a)$ has its second-to-last prefix decoding to a strict prefix of $x_1^n$, and every element of $\bar\gC_{n+1}(a)$ has its second-to-last prefix decoding to exactly $x_1^n$, a strict prefix of $x_1^na$; both are valid and decode to strings beginning with $x_1^na$. The two sets are disjoint because their second-to-last prefixes decode to strings of different lengths. Equation~\ref{eq:nextbyte} then follows from Lemma~\ref{lem:btr} applied to $x_1^na$ and $x_1^n$, using $P(\vt^\star\oplus\tau)=P(\vt^\star)P(\tau\mid\vt^\star)$ and $P(x_1^n)=\sum_{b}P(x_1^nb)$. \hfill$\square$

\section{Algorithmic details}
\label{app:algorithms}

Algorithm~\ref{alg:sampling} gives the complete byte-level sampling step corresponding to Figure~\ref{fig:sampling}. Starting from the prompt, $\cover(x_1^n)$ is computed by Algorithm~\ref{alg:cover}; each subsequent step reuses the returned cover set.

\begin{algorithm}[h]
\caption{Exact next-byte prediction and sampling from a tokenized LM}
\label{alg:sampling}
\begin{algorithmic}[1]
\Require byte string $x_1^n$, $\gC=\cover(x_1^n)$ with probabilities $\{P(\vt)\}_{\vt\in\gC}$, tokenized LM $P$, tokenizer $\enc$
\State $S(a)\gets 0$ and $\gC'(a)\gets\emptyset$ for all $a\in\gA$
\For{each $\vt\in\gC$ with $|\dec(\vt)|>n$} \Comment{step 2: regroup encodings extending beyond $n$}
  \State $a\gets\dec(\vt)_{n+1}$; \ $S(a)\gets S(a)+P(\vt)$; \ $\gC'(a)\gets\gC'(a)\cup\{\vt\}$
\EndFor
\State $\vt^\star\gets\enc(x_1^n)$ \Comment{the unique element of $\gC$ with $|\dec(\vt^\star)|=n$}
\State compute $P(\cdot\mid\vt^\star)$ with one forward pass \Comment{step 3: single LM call}
\For{each $\tau\in\gV$}
  \If{$\enc(\dec(\vt^\star\oplus\tau))=\vt^\star\oplus\tau$} \Comment{accept valid continuations only}
    \State $a\gets\tau_1$; \ $S(a)\gets S(a)+P(\vt^\star)P(\tau\mid\vt^\star)$; \ $\gC'(a)\gets\gC'(a)\cup\{\vt^\star\oplus\tau\}$
  \EndIf
\EndFor
\State $P(a\mid x_1^n)\gets S(a)/\sum_{b\in\gA}S(b)$ for all $a\in\gA$ \Comment{step 4: Equation~\ref{eq:nextbyte}}
\State sample $x_{n+1}\sim P(\cdot\mid x_1^n)$, or take the arg max for greedy decoding \Comment{step 5}
\State \Return $x_{n+1}$ and $\cover(x_1^{n+1})=\gC'(x_{n+1})$ with probabilities $\{P(\vt)\}$
\end{algorithmic}
\end{algorithm}

\textbf{Practical remarks.}
(i) All probabilities are maintained in log space, and sums are computed with log-sum-exp.
(ii) The validity check in Algorithm~\ref{alg:sampling} depends only on the tokenizer. Tokens whose probability is negligible do not affect Equation~\ref{eq:nextbyte} and can be skipped, which keeps the check inexpensive.
(iii) Ideally, $\sum_bS(b)=P(x_1^n)$; for a trained model, the normalization in step 4 discards the mass assigned to invalid encodings.
(iv) Special tokens (e.g., FIM sentinels and beginning-of-sequence tokens) are kept fixed as part of the context; the cover search operates only on the text after the last special token. An end-of-sequence token is mapped to a dedicated symbol appended to $\gA$.
(v) For ensembles, each member runs Algorithm~\ref{alg:sampling} independently up to step 4; the member distributions are averaged, a single byte is selected, and each member keeps its own $\gC'(x_{n+1})$.
(vi) For byte-level beam search, each hypothesis carries its own cover set; hypotheses are extended by all bytes with non-zero probability and ranked by their cumulative log-probability $\log P(x_1^n)$.

\section{Relation to token healing and token alignment}
\label{app:heuristics}
Token healing removes the last prompt token, i.e., conditions on $\vp=\enc(x_1^{n-|t_k|})$ where $t_k$ is the last token of $\enc(x_1^n)$, and restricts the next token to those having the removed bytes as a prefix. In the notation of Proposition~\ref{prop:structure}, this corresponds to retaining a single candidate boundary $j$ and discarding all others; furthermore, the retained boundary is chosen by the tokenizer rather than weighted by the model probability $P(\enc(x_1^j))$, and invalid continuations are not excluded. Token alignment extends this idea by backing up several tokens and constraining the generated tokens to match the removed characters, which again considers a subset of the cover encodings without weighting them according to Lemma~\ref{lem:btr}. Our method sums over every candidate boundary $j\in\{n-L,\dots,n-1\}$ with exact weights, as well as over all valid continuations of $\enc(x_1^n)$ during generation (Proposition~\ref{prop:recursion}), and therefore recovers the correct conditional distribution under the ideal model. Empirically, this difference translates into the gains over both heuristics in SPM mode reported in Figure~\ref{fig:fim}.
%%%%%%%%% Appendix %%%%%%%%%

\end{document}
